% Updated 4 November 2026 from docs/divergences.md and docs/chapter1--3
% at commit 2ed32808; checked against the book and the original erratum scan.
\documentclass[11pt]{article}
\usepackage[a4paper,margin=22mm]{geometry}
\usepackage{amsmath,amssymb,mathrsfs,longtable,array,bm}
\usepackage[hidelinks]{hyperref}
\usepackage{xurl}
\newcommand{\R}{\mathbf R}
\newcommand{\N}{\mathbf N}
\newcommand{\Leb}{\mathscr L}
\newcommand{\Haus}{\mathscr H}
\newcommand{\Mink}{\mathscr M}
\newcommand{\spt}{\operatorname{spt}}
\newcommand{\dmn}{\operatorname{dmn}}
\newcommand{\Clos}{\operatorname{Clos}}
\newcommand{\Hom}{\operatorname{Hom}}
\newcommand{\card}{\operatorname{card}}
\newcommand{\diam}{\operatorname{diam}}
\newcommand{\im}{\operatorname{im}}
\newcommand{\U}{\mathbf U}
\newcommand{\B}{\mathbf B}
\newcommand{\M}{\mathbf M}
\newcommand{\Lsp}{\mathbf L}
\newcommand{\Lt}{\widetilde{\mathbf L}}
\DeclareUrlCommand\rocq{\def\UrlFont{\footnotesize\ttfamily}}
\newcommand{\entry}[3]{#1&#2&#3\\[5pt]}
\newenvironment{errors}{%
  \begin{longtable}{@{}>{\raggedright\arraybackslash}p{19mm}
    >{\raggedright\arraybackslash}p{66mm}
    >{\raggedright\arraybackslash}p{72mm}@{}}
  \textbf{Item}&\textbf{Printed statement}&\textbf{Correction}\\\hline\endhead
}{\end{longtable}}
\newenvironment{confirmed}{%
  \begin{longtable}{@{}>{\raggedright\arraybackslash}p{19mm}
    >{\raggedright\arraybackslash}p{27mm}
    >{\raggedright\arraybackslash}p{111mm}@{}}
  \textbf{Item}&\textbf{Page, line}&\textbf{Correction adopted}\\\hline\endhead
}{\end{longtable}}
\title{\textit{Geometric Measure Theory}\\Errors detected by the formalisation}
\author{Compiled from the Rocq formalisation \texttt{federer-rocq}}
\date{4 November 2026}
\begin{document}
\maketitle

\noindent\textbf{Scope and sources (editorial).}
This list collects the statements of Herbert Federer, \textit{Geometric
Measure Theory} (1969), that the formalisation in this repository refutes as
printed, together with the corrections proved there. It also records
endpoint qualifications that depend on the formalisation's conventions.
The coverage is the formalised part of chapters 1 and 2, all of subsection
3.1, and the validated work through 3.2.48.
It is compiled from the divergence register \texttt{docs/divergences.md}
(rows of kind \emph{statement correction}) and from the section accounts in
\texttt{docs/chapter1}--\texttt{docs/chapter3}, at commit \texttt{2ed32808}.
The qualifications of 3.2.42--3.2.44 are registered there as conventions,
rather than unconditional refutations of the book.
Every entry was checked against the printed text of the Springer
reprint of the 1969 edition, which does not incorporate the 1987
corrections. Items are Federer's numbered paragraphs, and pages are printed
pages. The printed statements are paraphrased closely; displayed formulas
follow the book. Each entry names the Rocq results that prove the
counterexample and the corrected statement. Arguments that are not
formalised are marked as such.

Coverage does not imply completion of the deliberately deferred 1.1.3,
2.10.46, the second paragraph of 3.2.40 (Kneser's example), or the partial
3.2.45. The book examples excluded from the requested completion of
2.8.8--2.8.9 remain excluded. The supplements to 1.1.2, 1.1.4,
2.8.8--2.8.9 and 2.10.2 introduce no additional book errors to this list.

The check against the book removed one entry of the register. In 2.10.27,
p.~191, the book prints $f\,\B[(u,v),r]=\B(u,r)$, where the square brackets
group nested parentheses and $\B(u,r)$ is the closed ball (2.8.1). The
identity is correct; the register formerly read $\B(u,r)$ as an open
ball.

Federer's \textit{List of corrections} of May 1987 is transcribed in
\texttt{federer-erratum.tex}. Sections~\ref{s:false}--\ref{s:proofs}
contain local errors or convention qualifications absent from that list.
Section~\ref{s:erratum} records an error in that list and an editorial
normalization issue in another replacement.
Section~\ref{s:confirmed} lists the corrections of that list that the
formalisation adopts. The register formerly filed the entries for 2.7.14
and 2.7.15 as local corrections, but Federer's list already contains them (book
p.~131 line~10, p.~132 lines 22 and~25, p.~133 lines 3 and~4); they are
therefore listed in section~\ref{s:confirmed}.

Federer's notation is used: $\B(a,r)$ and $\U(a,r)$ are the closed and open
balls, $A\sim B$ is the set difference, and $\bigodot_m V$ is the $m$-th
symmetric power. The formalisation works over a field of characteristic
other than two where chapter 1 needs no order; the entries below are stated
over~$\R$.

\section{Statements that are false}\label{s:false}
\begin{errors}
\entry{1.4.4, p.~20}{For every associative algebra $W$ with unit, the
identity of $\Hom(V,W)$ extends to a homomorphism of the exterior algebra
of $\Hom(V,W)$ over $\R$ into the algebra of alternating $W$-valued forms.}
{False for noncommutative $W$: such an extension forces all values of a
one-form to commute. Take $V=\R^2$, $W$ the free associative algebra on two
letters $X,Y$, and $\alpha(e_1)=X$, $\alpha(e_2)=Y$; then
$(\alpha\wedge\alpha)(e_1,e_2)=XY-YX\neq0$. \emph{Correction:} require $W$
to be commutative; the homomorphism then exists and is unique.\newline
\rocq{noncommutative_exterior_extension_impossible}, \rocq{Omega_W},
\rocq{Omega_W_unique}}
\entry{1.4.4, p.~20}{The homomorphism of the preceding entry is surjective
whenever $\dim V<\infty$.}
{False for every $V$ unless $W=\R\cdot1_W$: in degree zero the map is
$c\mapsto c\cdot1_W$. For $W=\R\times\R$ the constant zero-form $(1,0)$ has
no preimage. \emph{Correction:} delete the assertion. In finite dimension
the isomorphism holds for the exterior algebra over a commutative $W$ of
the $W$-module $\Hom(V,W)$, as at the end of the same parenthesis.\newline
\rocq{coefficient_pair_not_surjective},
\rocq{Omega_W_surjective_requires_scalar_coefficients},
\rocq{Omega_relative_finite_bijective}}
\entry{2.4.12, p.~86}{Corollary of the Hahn--Banach theorem: for a
seminorm $\sigma$ on $V$ with values in $\{t:0\le t\le\infty\}$ and
$v\in V$ there is a linear $\alpha:V\to\R$ with $\alpha\le\sigma$ and
$\alpha(v)=\sigma(v)$.}
{False when $\sigma(v)=\infty$, since $\alpha$ is real valued. On $V=\R$
take $\sigma(0)=0$ and $\sigma(x)=\infty$ for $x\neq0$, and $v=1$.
\emph{Correction:} assume $\sigma(v)<\infty$.\newline
\rocq{infinite_seminorm_has_no_norming_functional},
\rocq{extended_seminorm_norming_functional}}
\entry{2.5.18(2), p.~112}{For $f,g,h:\R\to\R$ with $f,g$ continuous and
$g$ increasing, $\int_a^b(f\circ g)\,d(h\circ g)=\int_{g(a)}^{g(b)}f\,dh$;
``therefore $g_\#$ maps the measures induced by $h\circ g$ onto the measures
induced by~$h$''.}
{False unless $\im g=\R$. For $g(x)=x/(1+|x|)$ and $h(x)=x$, the image under
$g$ of the measure induced by $h\circ g$ vanishes on $\{t:2\le t\le3\}$,
whereas the measure induced by $h$ is $\Leb^1$.
The integral identity is unaffected. \emph{Correction:} assume
$\im g=\R$ in the final sentence.\newline
\rocq{stieltjes_image_unrestricted_false},
\rocq{stieltjes_image_requires_onto},
\rocq{stieltjes_image_induced_measures}}
\entry{2.6.5, p.~119}{The product of measures is associative, and
$\alpha_1\times\cdots\times\alpha_n$ does not depend on the association.}
{False for arbitrary measures. Let $\alpha=\beta=\Leb^1$ and let $\gamma$ be
a measure over a nonempty $Z$ with $\gamma(A)=\infty$ for every nonempty
$A\subset Z$. The product of the diagonal of $[0,1]$ and $Z$ has measure $0$
in $(\alpha\times\beta)\times\gamma$ and $\infty$ in
$\alpha\times(\beta\times\gamma)$. \emph{Correction:} assume that each
factor space is the union of countably many measurable sets of finite
measure (for the triple product, the first and third factors suffice).\newline
\rocq{lebesgue_product_not_associative},
\rocq{product_measure_associative_countable},
\rocq{finite_product_association_invariance}}
\entry{2.7.16(4), p.~138}{Weil's condition is violated by every
$t\in GL(n,\R)^W$ with $\det t=1$ and $\det(t|W)\neq1$.}
{The condition reads $\chi(t)\,|\det(t|W)|^{-n}\,|\det t|^m=1$, so it holds
when $\det(t|W)=-1$. \emph{Correction:} replace $\det(t|W)\neq1$ by
$|\det(t|W)|\neq1$. For $0<m<n$ the map acting as $2^{n-m}$ on $W$ and as
$2^{-m}$ on a complement has $\det t=1$ and $|\det(t|W)|=2^{m(n-m)}$, so
the conclusion stands.\newline
\rocq{grassmann_no_covariant_integral},
\rocq{stabiliser_scaling_exists}}
\entry{2.10.1, p.~170}{Alternate description: $\psi(A)$ is the infimum of
all $t$ such that every open covering $H$ of $A$ has a countable subfamily
$G$ of $F$ covering $A$, with each member in some member of $H$ and
$\sum_{S\in G}\zeta(S)<t$. ``For any topological space $X$, it remains
clear that $\psi$ measures $X$.''}
{False: the formula need not be monotone. Let $X=\{p,q\}$ with open sets
$\varnothing$, $\{q\}$, $X$, let $F=\{X\}$ and $\zeta(X)=1$. The value is
$\infty$ on $\{q\}$ and $1$ on $X$. \emph{Correction:} quantify over open
covers of the whole space $X$. The result is a measure, every continuous map
into a metric space is measurable, and it agrees with the printed formula
when $X$ is a metric space.\newline
\rocq{printed_topological_refinement_not_measure},
\rocq{topological_caratheodory_is_measure},
\rocq{topological_caratheodory_continuous_measurable},
\rocq{topological_caratheodory_metric_agreement}}
\entry{2.10.17(2), p.~179}{If $\mu$ is regular, $A\subset X$, $t>0$ and
$0\le(C)\limsup_{S\to x}\mu(A\cap S)/\zeta(S)<t$ whenever $x\in A$, then
$\mu(A)\le t\,\psi(A)$.}
{False. By 2.1.1 the quotient is undefined when $\zeta(S)=0$, and by
2.8.16 (p.~151) such $S$ are omitted from $(C)\limsup$. On $\R$ let every set belong to $F$,
$\zeta(S)=0$ for $\card S\le1$ and $\zeta(S)=1$ otherwise, $\mu$ the unit
mass at $0$, $A=\{0\}$. The upper density at $0$ is $1<2$, but
$\mu(A)=1>2\,\psi(A)=0$. \emph{Correction:} in (2) and (3) set $0/0=0$ and
$a/0=\infty$ for $a>0$; the density at $0$ is then $\infty$, and (2) and a
variant of (3) are proved. The original (3) is not refuted.\newline
\rocq{printed_caratheodory_density_bound_false},
\rocq{caratheodory_density_divide}, \rocq{federer_2_10_17_2},
\rocq{federer_2_10_17_3}}
\entry{3.1.17, p.~229}{``The converse of Theorem 3.1.16 is clearly also
true'': if for each $\varepsilon>0$ some map $g:\R^m\to\R^n$ of class $1$
satisfies $\Leb^m(A\sim\{x:f(x)=g(x)\})<\varepsilon$, then
$\operatorname{ap}\limsup_{x\to a}|f(x)-f(a)|/|x-a|<\infty$ for $\Leb^m$
almost all $a$ in~$A$.}
{False for nonmeasurable $A$: a constant map on a nonmeasurable subset of
$\R^m$ (2.2.4) has a smooth extension. Finite approximate upper growth
almost everywhere forces $A$ to be $\Leb^m$ measurable (via 3.1.8). The
counterexample uses the definition of 2.9.12 as corrected in Federer's list
(book p.~159 line~7), under which points outside $A$ count against the
density; under the original wording of p.~159 the example satisfies the
conclusion. \emph{Correction:} for $m>0$, finite approximate upper growth at almost all
points of $A$ is equivalent to $\Leb^m$ measurability of $A$ together with
the approximation property.\newline
\rocq{nonmeasurable_domain_converse_counterexample},
\rocq{approximate_growth_ae_domain_measurable}, \rocq{federer_3_1_17}}
\entry{3.2.27, p.~262}{If $W$ is an $(\Haus^m,m)$ rectifiable subset of
$\R^n$, $m\le n$ and $a_\lambda=\int N(\mathbf p_\lambda|W,y)\,d\Leb^my$ for
$\lambda\in\Lambda(n,m)$, then
$[\sum_\lambda(a_\lambda)^2]^{1/2}\le\Haus^m(W)\le\sum_\lambda a_\lambda$.}
{The upper bound is false without measurability of $W$ (3.2.14 does not
require it), even when the integrals exist. Take a
nonmeasurable $V\subset[0,1]$, put $V$ on the segment at height $0$ and
$[0,1]\sim V$ on the segment at height $1$. The coordinate integrals are
$1$ and $0$, but $\Haus^1(W)>1$. \emph{Correction:} assume $W$ is
$\Haus^m$ measurable for the upper bound. The lower bound holds as printed
whenever the integrals exist.\newline
\rocq{rectifiable_coordinate_upper_requires_measurability},
\rocq{rectifiable_coordinate_projection_lower}, \rocq{federer_3_2_27}}
\entry{3.2.37, p.~274}{The upper and lower Minkowski contents
$\Mink^{*m}$ and $\Mink_*^m$ are not measures; the preceding definition
permits every integer $0\le m\le n$.}
{False for $m=0$. Both contents are exactly counting measure on every
subset of $\R^n$, including infinite sets and the positive atom in
$\R^0$. \emph{Correction:} restrict the non-measure assertion to
$1\le m\le n$. For these positive degrees, finite sets have zero content
while an actual countable dense set has infinite content, so countable
additivity fails.\newline
\rocq{zero_minkowski_contents_counting},
\rocq{zero_minkowski_contents_are_measures},
\rocq{positive_minkowski_contents_not_measures}}
\entry{3.2.46, p.~281}{On an arbitrary open $A\subset\R^n$ with a
continuous field of inner products, the infimum of the lengths of
Lipschitzian curves joining $x,y\in A$ is a metric inducing the usual
topology of $A$.}
{The distance need not be finite. Take $A=\R\sim\{0\}$ with the ordinary
inner product. Every continuous curve from $-1$ to $1$ meets $0$, so
there is no admissible curve in $A$ and $\rho(-1,1)=\inf\varnothing=\infty$.
\emph{Correction:} allow an extended distance on $A$, or localize the
finite-metric assertion to a connected component. The formalisation
retains the original extended $\rho$; $\min\{1,\rho\}$ is a finite carrier
used for topology, with agreement of the original cutoffs below one.
The measure and transformation statements retain the original distance.\newline
\rocq{riemannian_punctured_line_infinite_distance},
\rocq{riemannian_distance_no_curve}}
\end{errors}

\section{Statements that fail in degenerate cases}\label{s:degenerate}
\noindent The literal definitions make the empty product $1$ the only simple
$0$-vector. Several assertions of chapter 1 therefore fail in degree zero.
Other assertions fail for zero spaces or empty sets.
\begin{errors}
\entry{1.4.4, p.~20}{The homomorphism of the first entry of
section~\ref{s:false} is not injective unless $\dim W=1$.}
{False for $V=0$ and $W=\R\times\R$: the map is $c\mapsto(c,c)$, which is
injective, although $\dim W=2$. \emph{Correction:} assume $V\neq0$ (and $W$
commutative). The corrected assertion is argued in \texttt{docs} by an
independent pair of one-forms, but is not formalised.\newline
\rocq{zero_domain_Omega_pair_injective}}
\entry{1.4.5, p.~20}{The ordinary tensor product of the exterior algebras of
$V$ and $W$ is neither commutative nor anticommutative (unless $V\cong\R$
or $W\cong\R$).}
{False for $V=W=0$: the tensor product is commutative, and neither space is
isomorphic to $\R$. Also false (not formalised) for $V=0$ and $\dim W\ge2$:
the product is then isomorphic to $\bigwedge_*W$, which is anticommutative
in the sense of 1.2.1. \emph{Correction (formalised):} the assertion holds
provided each of $V$ and $W$ contains two linearly independent vectors; no
finite dimension is needed. \emph{Sharp form (not formalised):} replace the
parenthesis by ``(unless $V=0$, or $W=0$, or $V\cong\R\cong W$)''; for
$V\cong\R$ and $\dim W\ge2$ the product is neither.\newline
\rocq{zero_ordinary_tensor_commutative},
\rocq{zero_space_not_isomorphic_scalar}, \rocq{ordinary_tensor_neither}}
\entry{1.6.1, p.~24}{A nonzero $m$-vector $\xi$ is simple if and only if
its associated subspace has dimension~$m$.}
{The reverse implication fails for $m=0$: the scalar $2$ has associated
subspace $\{0\}$ but is not simple. \emph{Correction:} assume $m>0$; for
$m=0$ only the forward implication holds.\newline
\rocq{simple_dimension_zero_counterexample}, \rocq{simple_iff_dimension}}
\entry{1.6.3, p.~24}{For $0\neq\phi\in\bigwedge^mV$ the associated
subspace $\bigwedge^1V\cap\{\alpha:\phi\wedge\alpha=0\}$ has dimension
$\le m$, with equality if and only if $\phi$ is simple.}
{Fails for $m=0$, with the covector $2$. \emph{Correction:} assume
$m>0$.\newline
\rocq{covector_dimension_zero_counterexample},
\rocq{simple_covector_iff_dimension}}
\entry{1.6.3, p.~24}{The isomorphisms $\mathbf D_p$ and $\mathbf D^p$ of
1.5.2 preserve simplicity.}
{Fails when the output degree is $0$, that is $p=n$: for $\dim V=1$,
$\mathbf D_1$ maps twice a basis vector, and $\mathbf D^1$ maps twice the
dual covector, to the scalar $2$. \emph{Correction:} assume $p<n$.\newline
\rocq{D_p_zero_degree_counterexample},
\rocq{D_upper_zero_degree_counterexample}, \rocq{D_p_preserves_simple}}
\entry{1.7.5, p.~32}{Equality holds in $(v_1\wedge\cdots\wedge v_m)\bullet
(v_1\wedge\cdots\wedge v_m)\le\prod_{i=1}^mv_i\bullet v_i$ if and only if $v_1,\ldots,v_m$ are mutually
orthogonal.}
{Fails for $(0,u,u)$ with $|u|=1$. \emph{Correction:} equality holds if and
only if some $v_i=0$ or $v_1,\ldots,v_m$ are mutually orthogonal.\newline
\rocq{hadamard_zero_factor_counterexample},
\rocq{exterior_product_norm_equality}}
\entry{1.8.1, p.~38}{The comass $\|\phi\|=\sup\{\langle\xi,\phi\rangle:
\xi\in\bigwedge_mV,\ \xi$ is simple, $|\xi|\le1\}$ satisfies
$|\phi|\ge\|\phi\|\ge\binom{\dim V}m^{-1/2}|\phi|$, and comass and mass are
dual norms.}
{Fails for $m=0$: the only simple $0$-vector is $1$, so $\|c\|=c$, and
$\|-1\|=-1$. \emph{Correction:} state comass,
mass and the norm assertions for $m>0$.\newline
\rocq{literal_zeroth_comass}, \rocq{literal_zeroth_comass_negative}}
\entry{1.8.1, p.~38}{The mass unit ball $C$ consists of all sums
$\sum_{i=1}^Nc_i\xi_i$ with $\xi_i\in S=\bigwedge_mV\cap\{\xi:\xi$ is simple
and $|\xi|\le1\}$, $c_i>0$, $\sum_{i=1}^Nc_i=1$ and
$N\le\dim\bigwedge_mV=\binom{\dim V}m$.}
{Fails when $\bigwedge_mV=0$, that is $m>\dim V$: then $N=0$ and the
coefficients cannot sum to $1$. \emph{Correction:} assume $m\le\dim V$. The
convex hull equality and the optimal decomposition hold for all $m$.\newline
\rocq{zero_dimensional_positive_convex_impossible},
\rocq{mass_positive_convex_representation}}
\entry{1.10.5, p.~48}{The mass $\|\xi\|$ of $\xi\in\bigodot_mV$ equals the
infimum of the sums $\sum_{j=1}^N|v_{1,j}|\cdots|v_{m,j}|$ over
representations $\xi=\sum_{j=1}^Nv_{1,j}\odot\cdots\odot v_{m,j}$; if
$\dim V<\infty$ one can take $N\le\dim\bigodot_mV$.}
{Fails for $m=0$: every empty product is $1$, so no finite sum of empty
products is $-1$, and the infimum for $\xi=-1$ is $\inf\varnothing=\infty$. \emph{Correction:} assume $m>0$; in degree zero the mass
of $c$ is $|c|$.\newline
\rocq{symmetric_negative_unit_not_empty_sum},
\rocq{symmetric_degree_zero_mass_specification}}
\entry{2.2.15, p.~71}{For an arbitrary separable metric space $Y$, every
Borel function mapping any metric space into $Y$ is a Baire function if and
only if for every finite subset $Z$ of $Y$ and every $\varepsilon>0$, $Z$ is
contained in the $\varepsilon$ neighbourhood of the image of some
continuous map of $\R$ into~$Y$.}
{Fails for $Y=\varnothing$: every map into $\varnothing$ is a Baire map,
whereas the condition for $Z=\varnothing$ demands a curve $\R\to
\varnothing$. \emph{Correction:} replace the condition by \emph{$Y$ is
empty or} the printed condition.\newline
\rocq{empty_target_borel_baire_counterexample},
\rocq{borel_baire_target_criterion_with_empty}}
\entry{2.3.2, p.~73}{A $\phi$ measurable function $f$, whose domain
contains $\phi$ almost all of $X$, can always be extended to a $\phi$ almost
equal function with domain~$X$.}
{Fails for $Y=\varnothing$: the empty function from a one-point space of
measure zero into $\varnothing$ has no total extension.
\emph{Correction:} assume $Y\neq\varnothing$.\newline
\rocq{empty_target_extension_counterexample},
\rocq{measurable_function_total_extension}}
\entry{2.3.6, p.~77}{Corollary of Lusin's theorem: in case $X$ is
countably $\phi$ measurable, $f$ is $\phi$ almost equal to some Borel function
mapping $X$ into~$Y$.}
{Fails for $Y=\varnothing$ and a nonempty $X$ of measure zero.
\emph{Correction:} assume $Y\neq\varnothing$ or $X=\varnothing$.\newline
\rocq{borel_version_empty_target_counterexample},
\rocq{lusin_borel_version_metric}, \rocq{lusin_borel_version_radon}}
\entry{2.7.3, p.~123}{The sets $\bigcup_gA_g(V)\times A_g(V)$, for nonempty
open $V$, form a basis of an invariant uniformity.}
{Fails for $X=\varnothing$: the family is empty, whereas a uniformity
contains $X\times X$. \emph{Correction:} assume $X\neq\varnothing$. The
equivalence in the first assertion holds also for $X=\varnothing$.\newline
\rocq{empty_space_orbit_family_not_basis_example},
\rocq{uniform_action_orbit_basis}}
\entry{2.7.4, p.~123}{There exists a nonzero monotone invariant integral on
$\mathscr K(X)$.}
{Fails for $X=\varnothing$: then $\mathscr K(X)=\{0\}$, although all other
hypotheses hold. \emph{Correction:} assume $X\neq\varnothing$.\newline
\rocq{invariant_integral_empty_space_example},
\rocq{uniform_action_invariant_integral_exists}}
\entry{2.10.28(4), p.~192}{If $J\in F$, $2\le n\in\mathbf Z$,
$h(s)=\zeta(J)/n$ and $t=(\diam J-s)/(n-1)$, then $0<s<\diam J/n$, $t>s$,
$\Phi(J,n)$ is special for $J$, and an estimate holds.}
{Fails when $J$ is a point: then $\zeta(J)=h(0)=0$, hence $s=0$ and
$t=0$.
\emph{Correction:} assume $\diam J>0$.\newline
\rocq{zero_parent_cantor_split_impossible},
\rocq{is_equal_cost_interval_family}}
\entry{2.10.38, pp.~198--199}{Every compact $S\subset\R^n$ has points
$a_i,b_i\in S$ and an enclosing box with decreasing side lengths $c_i$ such
that $|(a_1-b_1)\wedge\cdots\wedge(a_m-b_m)|=c_1\cdots c_m$ for
$m=1,\ldots,n$.}
{Fails for $S=\varnothing$ and $n>0$: no points $a_1,b_1$ exist.
\emph{Correction:} assume $S\neq\varnothing$.\newline
\rocq{empty_compact_enclosing_box_endpoint_counterexample},
\rocq{federer_2_10_38}}
\entry{2.10.44, p.~202}{For $S=\{(1,-1),(-1,1),(1,1)\}$, the supremum norm
$\mu$ and the Euclidean norm $\nu$, $\mu(u-v)=2=\nu[f(u)-f(v)]$ for
$u,v\in S$.}
{Each point has distance $0$ from itself. \emph{Correction:} distinct
points have distance~$2$.\newline
\rocq{kirszbraun_printed_two_needs_distinct},
\rocq{kirszbraun_corner_pair_distances}}
\entry{2.10.44, p.~202}{If $S$ is a subset of a metric space $X$, every
Lipschitzian $f:S\to\R$ has a Lipschitzian extension $g:X\to\R$ with
$\operatorname{Lip}(g)=\operatorname{Lip}(f)$; one may define
$g(\xi)=\inf\{f(x)+\operatorname{Lip}(f)\operatorname{dist}(x,\xi):x\in S\}$.}
{For $S=\varnothing$ and $X\neq\varnothing$ the infimum is $\infty$.
\emph{Correction:} assume $S\neq\varnothing$ for the formula; the
extension theorem holds for $S=\varnothing$ with $g=0$.\newline
\rocq{empty_source_no_finite_mcshane_envelope},
\rocq{federer_2_10_44_scalar_formula}, \rocq{federer_2_10_44_scalar}}
\entry{3.1.1, p.~209}{$\limsup_{x\to a}|f(x)-f(a)|/|x-a|
=\sup\{|\langle v,D f(a)\rangle|:|v|=1\}=\|D f(a)\|$, and the
corresponding formula for $\liminf$ and $\inf$.}
{Fails for $X=\{0\}$: there is no $x\neq a$ and no $v$ with $|v|=1$, so
by the conventions of 2.1.1 (p.~51) the first two expressions are
$\sup\varnothing=-\infty$, whereas $\|D f(a)\|=0$.
\emph{Correction:} assume $X\neq\{0\}$. The definition of the differential and the closed unit ball
formula of the norm include this case.\newline
\rocq{norm_ratio_zero_source_exception},
\rocq{norm_differential_sharp_limits}}
\entry{3.2.16, p.~253}{If $f$ is $(\phi,m)$ approximately differentiable
at $a$, with $g$ as in the definition, then $f$ is differentiable relative
to $\{x:f(x)=g(x)\}$ at~$a$.}
{The concept of 3.1.22 (p.~235) is introduced under the standing assumption
$a\in\Clos S$, which 3.2.16 does not ensure. For $\phi=0$ and the empty
function every point is a point of approximate differentiability, while
every agreement set is empty. Without the standing assumption the
definition of 3.1.22 is satisfied vacuously, so the defect is a missing
presupposition, not a false conclusion. \emph{Correction:} add
$a\in\Clos\{x:f(x)=g(x)\}$, which holds when $\phi\{a\}>0$.\newline
\rocq{measure_approximate_empty_zero},
\rocq{measure_approximate_empty_agreement_not_closure},
\rocq{differentiable_representative_relative_agreement}}
\end{errors}

\section{Convention-dependent endpoint qualifications}\label{s:conventions}
\noindent These entries concern the zero exponent in dimension one.
The formalisation uses $0^0=1$; the following are not claims that the book
explicitly assigns that value in these formulas. All cases $n\ge2$ are
proved as printed, including zero volumes and empty cores.
\begin{errors}
\entry{3.2.42, p.~278}{The outer and inner distance-strip volumes,
divided by $r>0$, are bounded below by
$n\alpha(n)^{1/n}v^{(n-1)/n}$, where $v$ is respectively the closure
volume or the volume of the inner core
$T_r=\{x:\operatorname{dist}(x,\R^n\sim S)\ge r\}$.}
{For $n=1$, the exponent is zero. With $0^0=1$ the inner formula fails
for $S=(0,1)$ and $r=2$: $T_2=\varnothing$, and the inner strip volume
divided by two is at most $1/2$, whereas the right side is
$\alpha(1)=2$. \emph{Qualification:} in dimension one retain a nonempty
source for the outer bound and a nonempty core for the inner bound.
The formalisation proves both such bounds, and both unrestricted bounds
for $n\ge2$.\newline
\rocq{boundary_layer_empty_core_line_exclusion_example},
\rocq{outer_boundary_layer_volume_nonempty},
\rocq{inner_boundary_layer_volume_nonempty_core},
\rocq{federer_3_2_42}}
\entry{3.2.43--44, p.~278}{The lower $(n-1)$-dimensional Minkowski
content of $\operatorname{Bdry}S$ is at least
$n\alpha(n)^{1/n}v^{(n-1)/n}$, where
$v=\Leb^n(\Clos S)<\infty$. For an $(n-1)$ rectifiable boundary,
3.2.44 replaces its lower content by its Hausdorff measure.}
{For $n=1$ and $S=\{0\}$, $v=0$ and the boundary lower degree-zero
content is one. Substituting $0^0=1$ would demand at least two.
\emph{Qualification:} retain positive closure volume in the line case;
all finite closure volumes, including zero, remain in every $n\ge2$.
The same qualification is retained in the rectifiable-boundary
corollary. No convention assigning a value to $0^0$ is inferred from
these printed inequalities.\newline
\rocq{isoperimetric_null_line_singleton_exclusion_example},
\rocq{positive_closure_volume_isoperimetric},
\rocq{federer_3_2_43}, \rocq{federer_3_2_44}}
\end{errors}

\section{Misprints in statements, formulas and the index}\label{s:misprints}
\begin{errors}
\entry{1.9.1, p.~41}{$\bigodot_*V=\bigoplus_{m=2}^\infty\bigodot_mV$.}
{Read $m\ge0$, as the next sentence (degrees zero and one) requires. The
sum for the ideal starts correctly at $2$.\newline
\rocq{symmetric_zero_ambient_empty_word}}
\entry{1.9.4, p.~43}{$\Upsilon(e^\alpha)=\sum_{p=0}^m\sum
\binom{\alpha_1}{\beta_1}\cdots\binom{\alpha_m}{\beta_m}e^\beta\otimes
e^{\alpha-\beta}$ for $\alpha\in\Xi(n,m)$, the inner sum over
$\alpha\ge\beta\in\Xi(n,p)$.}
{The multi-indices have $n$ entries; read
$\binom{\alpha_1}{\beta_1}\cdots\binom{\alpha_n}{\beta_n}$.\newline
\rocq{symmetric_diagonal_monomial}}
\entry{1.10.6, p.~49}{$\alpha(j)=\card\{i:s(i)=j\}$ for $j=1,\ldots,m$.}
{Read $j=1,\ldots,n$.\newline
\rocq{symmetric_prescribed_fibre_condition},
\rocq{symmetric_fibre_cardinality}}
\entry{2.7.16(7), p.~141}{$s[p^{-1}\{y\}]=(p\circ s^{-1})^{-1}\{y\}$ for
$s\in\mathbf O(n)$, $y\in\R^n$.}
{Read $y\in\R^{n-m}$, since $p$ maps $\R^n$ into $\R^{n-m}$.\newline
\rocq{affine_parametrisation_fibre}}
\entry{3.2.22, p.~260}{$\int_W\operatorname{ap}J_\mu f\,d\Haus^m=\int_Z
\Haus^{m-\mu}(W\cap f^{-1}\{z\})\,d\Haus^mz$.}
{Read $d\Haus^\mu z$, as in assertion (3) on p.~258 and in the preceding
integrals.\newline
\rocq{rectifiable_coarea_subset_formula}, \rocq{federer_3_2_22_3}}
\entry{Index, p.~673}{The entry for Gross measure points to 2.10.3,
p.~172.}
{Read 2.10.4(1), p.~172.\newline
\rocq{gross_measure}}
\end{errors}

\section{Errors in proofs}\label{s:proofs}
\noindent The statements of these items are correct; an intermediate
assertion of the printed proof is false.
\begin{errors}
\entry{2.10.24, p.~188}{Remark after the proof: the closed balls in a
boundedly compact metric space $X$ form a closed set with respect to the
Hausdorff distance.}
{False. In $X=[-1,-1/2]\cup[1/2,1]\cup\{-2,2\}\subset\R$ the set
$C=[-1,-1/2]\cup[1/2,1]$ is not a ball, but the balls with centre $1/2+h$
and radius $3/2-2h$ lie within Hausdorff distance $3h$ of $C$.
\emph{Correction:} the spherical equality is proved by Hausdorff
comparison instead.\newline
\rocq{closed_ball_hausdorff_closedness_counterexample},
\rocq{spherical_weighted_proper_equality}}
\entry{2.10.31, p.~196}{After the reflected-cap formula, $T_k\subset\B(a,s)$
for some $s<r$.}
{Here $r=\inf\{s:T\subset\B(0,s)$ for some $T\in C\}$ and $a$ lies on the
sphere $\operatorname{Bdry}\B(0,r)$, so the contradiction requires a ball
with centre $0$. The origin singleton lies in the centred unit ball and
misses its sphere, but lies in no ball of radius less than $1$ centred at a
point $a$ with $|a|=1$. \emph{Correction:} read $T_k\subset\B(0,s)$ for some
$s<r$.\newline
\rocq{shrinking_ball_wrong_center_counterexample},
\rocq{steiner_minimal_radius_contains_sphere}}
\entry{3.1.24, p.~241}{The norms of all terms of the coordinate series
total at most $\sum_{m=0}^\infty c^mm!\,(\sum_{i=1}^n|x_i-a_i|)^m/m!$;
the term of order $0$ is thus bounded by~$1$.}
{Condition (1) bounds only the differentials of positive order; the constant
analytic function $2$ violates the printed bound at the centre.
\emph{Correction:} keep $|f(a)|$ as a separate term. Absolute convergence
is unaffected.\newline
\rocq{analytic_printed_constant_term_bound_fails},
\rocq{analytic_coordinate_zero_degree_cost}}
\entry{3.2.29, p.~268}{$(\im F)\sim\bigcup\{f(U):U\in G\}
=F(\R^m\sim C)\cup f(C\sim P)$.}
{Only the inclusion $\subset$ holds: $F(\R^m\sim C)$ may meet
$\bigcup\{f(U):U\in G\}$. \emph{Correction:} replace $=$ by $\subset$,
which suffices for the subsequent estimate.\newline
\rocq{agreement_image_cover_remainder},
\rocq{lipschitz_small_manifold_remainder}, \rocq{federer_3_2_29}}
\entry{3.2.38, p.~275}{The first displayed Steiner sum in the bound
for $\rho(K)$ runs from $j=1$ to $n$.}
{It omits the constant intrinsic-volume term. For the actual zero affine
map on the unit line cube, $m=n=1$, the image is $\{0\}$ and its
normalized radius-$r$ tube volume is $2r>0$, while the only printed
$j=1$ coefficient is zero. \emph{Correction:} sum from $j=0$ to $n$.
The extra term is absorbed in the unspecified finite error constant;
the statement of the lemma is unchanged.\newline
\rocq{missing_steiner_constant_counterexample_exists},
\rocq{compact_minkowski_constant_map_correction_example},
\rocq{affine_coordinate_cube_steiner_estimate}}
\end{errors}

\section{An error and a normalization issue in the list of corrections
of May 1987}\label{s:erratum}
\begin{errors}
\entry{2.9.18(1), p.~163}{Erratum p.~3, book p.~163 line~11: in
$\psi(A)=\sup(\{0\}\cup\{c_j:A\cap S_j\neq\varnothing\})$ read
$A\cap S_j=\varnothing$.}
{With this replacement $\psi(\varnothing)=c_1=1$, so $\psi$ is not a
measure, and the computation $\psi\B(x,r)=c_{j+1}$ for
$c_{j+1}\le r\le jc_{j+1}$ fails. \emph{Correction:} keep the book's
condition $A\cap S_j\neq\varnothing$.\newline
\rocq{weighted_miss_capacity_empty_positive},
\rocq{weighted_miss_capacity_not_measure},
\rocq{weighted_hit_measure_is_measure}}
\end{errors}

\noindent\textbf{Normalization note for 3.2.38 (editorial).}
On erratum p.~5, book p.~275 line~2, the replacement is $r^{n-j}$.
The original scan confirms this exponent. In the displayed bound for
$\rho(K)$, however, the tube volume has already been divided by
$\alpha(n-m)r^{n-m}$ (definition on book p.~274). A raw Steiner term has
power $r^{n-j}$; its normalized power is therefore $r^{m-j}$, as printed
in the book. The formalisation keeps the full raw Steiner polynomial
and divides by the actual denominator, retaining $j=0$ as above; see
\rocq{affine_coordinate_cube_steiner_estimate}.
Thus the erratum's replacement cannot be applied literally to that
already normalized power. This is an algebraic consistency observation;
no separate Rocq counterexample to this erratum replacement is claimed.
The transcription \texttt{federer-erratum.tex} preserves the scan.

\section{Corrections of the list of May 1987 confirmed by the
formalisation}\label{s:confirmed}
\noindent The formalisation adopts the following corrections of Federer's
list. Pages and lines are those of the book. The table includes corrections
used in the formalised statements or arguments; purely stylistic and
bibliographical changes are omitted.
\begin{confirmed}
\entry{1.5.1}{22, 9}{In the associativity identity of interior
multiplication the multivector is $\zeta\in\bigwedge_tV$.}
\entry{1.7.8}{34, 29--30}{$\langle\xi,*\beta\rangle=(-1)^{p(n-p)}\langle
*\xi,\beta\rangle$ for $\xi\in\bigwedge_pV$, $\beta\in\bigwedge^{n-p}V$.}
\entry{1.7.9}{35, 3}{$(v\otimes w)\bullet(v'\otimes w')=(v\bullet
v')(w\bullet w')$.}
\entry{2.1.5}{57, 26}{$\phi(A)=\card(A)$ for $A\subset S$.}
\entry{2.2.1}{60, 4}{Both generation criteria assume $\varnothing\in S$;
for $S=\varnothing$ the generated family $\{\varnothing\}$ is not a Borel
family over a nonempty $X$.}
\entry{2.2.10}{66, 13 and 15}{A sequence $S_1,S_2,S_3,\ldots$ of closed
sets in the cylinders $\mathscr N_j$.}
\entry{2.2.10}{67, 2}{$\bigcup G$ and $\bigcup H$ are Borel separated.}
\entry{2.3.8}{78, 7}{$\mathbf A(\phi,Y)$ is a set of partial functions; the
vector space and convex hull statements concern its members with domain
$X$.}
\entry{2.4.12}{86, 26--27 and bottom; 87, 1}{$\phi_{(p)}$ is defined on
$\mathbf A(\phi,Y)$ and is a seminorm on $\widetilde{\mathbf A}(\phi,Y)$.}
\entry{2.4.12}{87, 18}{The functions with finite image are dense in
$\Lsp_p(\phi,Y)$ provided $p<\infty$ or $\dim Y<\infty$.}
\entry{2.4.12}{87, 32; 88, 3}{$\phi_{(1)}$ closed subspace; the subspace of
functions with finite image.}
\entry{2.4.16}{89, 19}{The tests range over $u\in\Lsp_q(\phi)$ with
$\phi_{(q)}(u)\le1$.}
\entry{2.5.1}{91, 28}{$\sup\{f,g\}\in L$ for all $f,g\in L$.}
\entry{2.5.7}{98, 10--11}{$\mu$ is defined on $\Lt_p(\phi)$, and the
supremum is over $\phi_{(p)}(f)\le1$.}
\entry{2.5.7}{99, 13, 14, 17}{The residual Daniell integral is taken on
$L^+$.}
\entry{2.5.8}{100, 3 and 5}{The Daniell integral is on $\Lt_\infty(\phi)$
and vanishes whenever $\phi_{(\infty)}(f)=0$; without this hypothesis
evaluation at a point under the zero measure is a counterexample.}
\entry{2.5.9}{100, 21--24}{$\Lt_1(\phi)$ in place of $\Lsp_1(\phi)$;
$L^+$ is dense in $\Lt_1(\phi)^+$.}
\entry{2.5.11}{102, 8 and 12}{The lattice is
$\Lt_1(\alpha_1\times\beta_2)\cap\Lt_1(\beta_1\times\alpha_2)$, and
$\Lt_1(\phi)=L$.}
\entry{2.6.4}{118, 20, 21, 24; 119, 3, 4, 6}{$\Lt_1$ in place of
$\Lsp_1$.}
\entry{2.6.4}{119, 19}{$\nu(h)=\int h\,d(\alpha\times\beta)$ for
$h\in\mathscr K(X\times Y)$.}
\entry{2.7.4}{124, 30; 125, 7--8}{``Whenever $0\neq\nu\in\mathscr K(X)^+$,
$(\spt\nu)\times(\spt\nu)\subset\beta$, $a\in\spt\nu$'', and
$\lambda\in S(\beta)$.}
\entry{2.7.14}{131, 10}{The measure associated with a \emph{nonzero}
monotone left (right) $\chi$ covariant integral is a Radon measure. For the
zero integral on the group of subsets of $\N$ under symmetric difference the
associated measure is not a Radon measure
(\rocq{zero_left_invariant_measure_not_radon}).}
\entry{2.7.14}{132, 22 and 25}{$f\in\Psi$ in place of $f\in\Phi$.}
\entry{2.7.14}{133, 3}{$U=(G\sim P)\cup(P\cap U)$.}
\entry{2.7.15}{133, 4}{The measure associated with a \emph{nonzero}
monotone $\chi$ covariant integral is a Radon measure
(\rocq{homogeneous_covariant_measure_radon}).}
\entry{2.7.16(1)}{134, 26; 135, 35}{$\Leb^n$ is the unique \ldots; the
exponent in the recursion for $\bm\alpha(n)$ is $(n-1)/2$.}
\entry{2.7.16(2)}{136, 29}{The left Haar integral of $GL(V)$ is $\int
u\,|\det|^{-n}\,d\phi$.}
\entry{2.7.16(5)}{138, 34}{The modular function of $G^W$ involves
$|\det[L(g)|W]|^{-1}$.}
\entry{2.7.18}{141, 25}{$\chi(g)\to\chi(h)$.}
\entry{2.8.2}{142, 33}{$G_0=\varnothing$.}
\entry{2.8.4}{143, 17}{$\varnothing\notin F$.}
\entry{2.8.17}{151, 24; 152, 10}{$\phi$ is positive on $V(X)$; cite
2.8.3.}
\entry{2.9.7}{155, 26}{$\mathbf D(\psi,\phi,V,x)\notin\R$.}
\entry{2.9.12}{159, 7}{$X\sim\{z:f(z)\le t\}$ has density $0$ at $x$.}
\entry{2.9.13}{159, 19}{$(\phi,V)$.}
\entry{2.9.15}{161, 11}{Cite 2.8.14.}
\entry{2.9.17}{162, 24--25}{$\B(x,2r)$.}
\entry{2.10.1}{170, 24 and 35}{Cite [F12].}
\entry{2.10.14}{177, 23}{The upper integral $\int^*$.}
\entry{2.10.21}{184, 8 and 10}{Nonempty compact convex subsets.}
\entry{2.10.28(3)}{192, 4}{$\card(H_0)=1$.}
\entry{2.10.28(4)}{193, 7}{Sum over $S\in\Phi(j,n)$ with $S\subset T$.}
\entry{2.10.29}{194, 8}{$\sigma_j\ge t_{j+1}\to0$.}
\entry{2.10.32}{197, 4}{$\mu\bm\alpha(m)^{-1}2^m$.}
\entry{3.1.1}{210, 26}{$u_{n+1}-u_n=g(u_{n-1})-g(u_n)$.}
\entry{3.1.4}{214, 24}{$\ge1/j$.}
\entry{3.1.9}{218, 12}{$\U(z,1/j)$.}
\entry{3.1.11}{219, 4}{$\{|\phi(v_1,\ldots,v_m)|:\ldots\}$.}
\entry{3.1.11}{221, 1}{$a+t(x-a)$.}
\entry{3.1.12}{224, 4}{$\alpha h(x)+(\beta+1)h(s)$.}
\entry{3.1.16}{228, 19}{All $a$ in $A$.}
\entry{3.1.20}{232, 11, 14--15}{Maps $f_i:W_i\to B$ of class $k$ with
$f_i(x)=x$ for $x\in B\cap W_i$; the blend maps into $T\cap B$.}
\entry{3.1.22}{235, 19}{$f(a)=\eta$ in case $a\in S$.}
\entry{3.2.2}{242, 23}{$\mathbf O(n,m)$.}
\entry{3.2.3}{244, 11}{$J_mf(x)$.}
\entry{3.2.6}{246, 3 and 13}{$f(a)\le f(b)$; $\Leb^1$ summable.}
\entry{3.2.10}{248, 21 and 25}{$\Leb^{m-n}$ and $\Haus^{m-n}$.}
\entry{3.2.11}{249, 21}{$\Leb^m(A)$.}
\entry{3.2.16}{253, 29}{$\delta_{i+1}<\delta_i/2$.}
\entry{3.2.16}{253, 32}{Differentiable relative to $T$ at $a$; $T\cup\{a\}$
in case $\phi\{a\}>0$.}
\entry{3.2.16}{254, 2}{``generalizes'' in place of ``coincides with''.}
\entry{3.2.21}{257, 26}{$\Leb^\mu(B)$.}
\entry{3.2.22}{258, 20 and 22; 259, 3--4}{The source remainder
$R=W\sim\bigcup_{i=1}^\infty\psi_i(K_i)$, with $\Haus^m(R)=0$ and
$\Haus^{m-\mu}$ null fibres.}
\entry{3.2.25}{261, 4}{$\Haus^m$.}
\entry{3.2.26}{261, 22 and 29; 262, 14}{Cite 2.10.5; $x-w$; cite 3.2.20.}
\entry{3.2.28(2)}{265, 3}{Delete ``because $\mathbf O(n)$ is compact,''.}
\entry{3.2.28(3)}{265, 12}{$\mathbf O(m,m-1)$.}
\entry{3.2.33}{270, 25 and 28}{Cite 3.1.24 for the analytic rank
alternative, and 3.1.18 for positivity of the complementary-dimensional
fibres.}
\entry{3.2.35}{271, 24}{Parametrize the convex boundary by
$x\mapsto g(x)^{-1}x$ on the unit sphere, where $g$ is the gauge.}
\entry{3.2.36}{273, 18}{Read ``$m$-th'' in the elementary symmetric
function description of the box coefficient.}
\entry{3.2.38}{275, 14}{Use the actual radius-$r$ tube
$\{z:\operatorname{dist}[z,f(C)]<r\}$ in the final volume quotient.}
\end{confirmed}
\end{document}
